
\documentclass{article}
\usepackage{colortbl}
\usepackage{makecell}
\usepackage{multirow}
\usepackage{supertabular}

\begin{document}

\newcounter{utterance}

\centering \large Interaction Transcript for game `wordle\_withcritic', experiment `medium\_frequency\_words\_no\_clue\_no\_critic', episode 15 with Ministral{-}3{-}14B{-}Reasoning{-}2512{-}nothink{-}t0.0.
\vspace{24pt}

{ \footnotesize  \setcounter{utterance}{1}
\setlength{\tabcolsep}{0pt}
\begin{supertabular}{c@{$\;$}|p{.15\linewidth}@{}p{.15\linewidth}p{.15\linewidth}p{.15\linewidth}p{.15\linewidth}p{.15\linewidth}}
    \# & $\;$A & \multicolumn{4}{c}{Game Master} & $\;\:$B\\
    \hline

    \theutterance \stepcounter{utterance}  
    & & \multicolumn{4}{p{0.6\linewidth}}{
        \cellcolor[rgb]{0.9,0.9,0.9}{
            \makecell[{{p{\linewidth}}}]{
                \texttt{\tiny{[P1$\langle$GM]}}
                \texttt{You are a language wizard who likes to guess words by using the given rules.} \\
\\ 
\texttt{Welcome to Wordle! You have six attempts to guess the target word, a valid English word of five lowercase letters (a{-}z). Please use the tags "explanation:" and "guess:" to provide a concise explanation for each guess.} \\
\\ 
\texttt{For instance, if your guess is "apple", your response should be} \\
\texttt{explanation: this is a common five{-}letter English word, and I am starting my guess with this word.} \\
\texttt{guess: apple} \\
\\ 
\texttt{After each guess, your answer will be validated, and you will receive feedback indicating which letters are correct (green), which letters are correct but in the wrong position (yellow), and which letters are incorrect (red). This feedback can be useful in determining which letters to include or exclude in your next guess.} \\
\\ 
\texttt{For example, the feedback for "apple" might be:} \\
\texttt{guess\_feedback: a<yellow> p<yellow> p<green> l<yellow> e<red>} \\
\\ 
\texttt{The explanation should contain details about how the guess\_feedback is used to arrive at a new guess.} \\
\\ 
\texttt{Do not generate any other text.} \\
\\ 
\texttt{Let's begin with your first guess.} \\
\\ 
\texttt{clue: passenger ship} \\
            }
        }
    }
    & \\ \\

    \theutterance \stepcounter{utterance}  
    & \multicolumn{4}{p{0.6\linewidth}}{
        \cellcolor[rgb]{0.9,0.9,0.9}{
            \makecell[{{p{\linewidth}}}]{
                \texttt{\tiny{[P1$\rangle$GM]}}
                \texttt{explanation: I'm starting with a common five{-}letter word that might contain letters relevant to "passenger ship." I'll use "ship" as a base and add a common prefix to form a valid word.} \\
\texttt{guess: ships} \\
            }
        }
    }
    & & \\ \\

    \theutterance \stepcounter{utterance}  
    & & \multicolumn{4}{p{0.6\linewidth}}{
        \cellcolor[rgb]{0.9,0.9,0.9}{
            \makecell[{{p{\linewidth}}}]{
                \texttt{\tiny{[GM$\rangle$P2]}}
                \\ 
\\ 
\texttt{clue: passenger ship} \\
\texttt{explanation: I'm starting with a common five{-}letter word that might contain letters relevant to "passenger ship." I'll use "ship" as a base and add a common prefix to form a valid word.} \\
\texttt{guess: ships} \\
            }
        }
    }
    & \\ \\

    \theutterance \stepcounter{utterance}  
    & & & \multicolumn{4}{p{0.6\linewidth}}{
        \cellcolor[rgb]{0.9,0.9,0.9}{
            \makecell[{{p{\linewidth}}}]{
                \texttt{\tiny{[GM$\langle$P2]}}
                \texttt{Your guess of **"ships"** is a good start, but let's refine it based on the clue **"passenger ship"** and the explanation provided.} \\
\\ 
\texttt{\#\#\# Analysis:} \\
\texttt{1. **Clue: "passenger ship"**} \\
\texttt{   {-} A "passenger ship" is a vessel designed to carry people (passengers) across water, often for travel or transportation.} \\
\texttt{   {-} Common terms related to this could include:} \\
\texttt{     {-} **Passenger** (but this is 8 letters, not 5)} \\
\texttt{     {-} **Ferry** (a type of passenger ship, but 5 letters)} \\
\texttt{     {-} **Cruise** (a type of passenger ship journey, but 6 letters)} \\
\texttt{     {-} **Liner** (a passenger ship, but 5 letters)} \\
\texttt{     {-} **Steam** (as in "steamship," but 5 letters)} \\
\texttt{     {-} **Ocean** (as in "ocean liner," but 6 letters)} \\
\texttt{     {-} **Vessel** (too long)} \\
\texttt{     {-} **Cabin** (part of a passenger ship, but 5 letters)} \\
\texttt{     {-} **Deck** (part of a passenger ship, but 4 letters)} \\
\texttt{     {-} **Port** (related to passenger ships, but 4 letters)} \\
\texttt{     {-} **Harbor** (too long)} \\
\texttt{     {-} **Yacht** (a type of passenger ship, but 5 letters)} \\
\\ 
\texttt{2. **Explanation: Starting with a common five{-}letter word, using "ship" as a base, and adding a common prefix.**} \\
\texttt{   {-} The base word is **"ship"** (4 letters).} \\
\texttt{   {-} Adding a common prefix (1 letter) to make a 5{-}letter word:} \\
\texttt{     {-} **A{-}** (e.g., "aship" – not a word)} \\
\texttt{     {-} **Be{-}** (e.g., "beship" – not a word)} \\
\texttt{     {-} **De{-}** (e.g., "deship" – not a word)} \\
\texttt{     {-} **En{-}** (e.g., "enship" – not a word)} \\
\texttt{     {-} **In{-}** (e.g., "inship" – not a word)} \\
\texttt{     {-} **On{-}** (e.g., "onship" – not a word)} \\
\texttt{     {-} **Re{-}** (e.g., "reship" – not a word)} \\
\texttt{     {-} **Un{-}** (e.g., "unship" – not a word)} \\
\texttt{     {-} **Pre{-}** (e.g., "preship"} \\
            }
        }
    }
     \\ \\

    \theutterance \stepcounter{utterance}  
    & & & \multicolumn{2}{p{0.3\linewidth}}{
        \cellcolor[rgb]{0.9,0.9,0.9}{
            \makecell[{{p{\linewidth}}}]{
                \texttt{\tiny{[GM$|$GM]}}
                \texttt{The response should always start with the keyword 'explanation:'} \\
            }
        }
    }
    & & \\ \\

    \theutterance \stepcounter{utterance}  
    & & & \multicolumn{2}{p{0.3\linewidth}}{
        \cellcolor[rgb]{0.9,0.9,0.9}{
            \makecell[{{p{\linewidth}}}]{
                \texttt{\tiny{[GM$|$GM]}}
                \texttt{game\_result = ABORT} \\
            }
        }
    }
    & & \\ \\

\end{supertabular}
}

\end{document}
